% METODOLOGÍA (capítulo IV con hipótesis, III sin hipótesis). Lo escribe tesis-metodologia.
% Las secciones cambian entre plan y tesis; \ifmodoplan elige cuál se imprime.

\section{Diseño de investigación}
% \subsection{Tipo}, \subsection{Nivel}, \subsection{Diseño} (esquema GE: O1 X O2),
% \subsection{Enfoque}, \subsection{Método}; cada uno definido con cita y aplicado.

\section{Descripción de la metodología}
% Cómo se ejecuta el diseño: grupo, cuándo se miden O1 y O2, indicadores e instrumentos.
% Separa el desarrollo del software de la comprobación de hipótesis.
\ifmodoplan
\subsection{Etapas del desarrollo del plan de tesis}
% Etapas de la investigación y del desarrollo del sistema (p. ej. Scrum: sprint 0 ... n).
\else
\subsection{Implementación del tema de investigación}
% Desarrollo del sistema por sprints: requerimientos, arquitectura, modelo de datos,
% integraciones (WhatsApp, OpenAI), capturas y fragmentos de código en lstlisting.
\subsection{Pruebas realizadas}
% Pruebas funcionales y cómo se tomaron las mediciones pre y post de cada indicador.
\fi

\ifconhipotesis
\section{Población y muestra}
% \subsection{Población} y \subsection{Muestra}: unidad de análisis, criterios de inclusión y
% exclusión, tamaño (fórmula si es probabilístico) y tipo de muestreo.

\ifmodoplan
\section{Técnicas e instrumentos de recolección de datos}
% Técnica -> instrumento por indicador, validez y confiabilidad previstas, plan de análisis.
\else
\section{Técnicas de recolección de datos}
% Técnica usada por indicador, con cita.

\section{Instrumentos de recolección de datos}
% Instrumento por indicador (fichas en el Anexo 2).
\subsection{Validez}
% Juicio de expertos (3): claridad, pertinencia, relevancia. Resultado en el Anexo 3.
\subsection{Confiabilidad}
% Datos del sistema: trazabilidad de registros. Cuestionarios: Alfa de Cronbach > 0,7.
\fi
\fi
